% Rewritten Section VII generated from the 20-seed results under output/.
% Required packages: graphicx, booktabs, amsmath.
% The paths below assume that the main LaTeX document is compiled from the
% repository root. Adjust the paths if this file is copied elsewhere.

\section{Experimental Setup, Results, and Discussion}
\label{sec:experiments}

This section evaluates the Unified Bayesian-Robust Credit-Aware
Scheduler (UBR-CA) with respect to execution cost, deadline compliance,
CPU-credit safety, migration overhead, parameter sensitivity, and
scalability. The study addresses the following research questions:
\emph{RQ1}: What execution-cost overhead or saving does UBR-CA produce
relative to conventional workflow and credit-aware schedulers?
\emph{RQ2}: Does robust credit-aware scheduling improve workflow
deadline compliance?
\emph{RQ3}: How effectively does UBR-CA prevent credit exhaustion under
increasing credit stress?
\emph{RQ4}: What are the individual contributions of Bayesian
uncertainty learning and proactive migration?
\emph{RQ5}: How sensitive is the method to its principal control
parameters, and does it remain practical for large workflows?

\subsection{Experimental Protocol and Workloads}
\label{subsec:protocol}

The numerical results reported in this section were obtained from
benchmark-shaped synthetic workflow DAGs with Montage, Epigenomics,
LIGO, and production-trace-like structural characteristics. The
workload generator preserves the principal differences among these
classes: Montage-like workflows have broad parallel stages,
Epigenomics-like workflows combine chains and branches, LIGO-like
workflows contain deeper dependency paths, and the production-like
class uses randomized acyclic precedence relations. Each default run
contains eight concurrently arriving workflows with approximately 90
tasks per workflow. A deadline is assigned as 1.55 times the
workflow's nominal critical-path duration.

The implementation also provides a CSV importer for replaying
preprocessed Alibaba Cluster Trace workflows. However, no processed
Alibaba trace was supplied to the reported run. Consequently, the
values below must be interpreted as synthetic benchmark results and
must not be described as Alibaba trace measurements. The trace importer
can regenerate the complete table and figure set when trace-backed
tasks, arrivals, utilization values, and dependencies are available.

Scalability experiments use the workflow-size categories in
Table~\ref{tab:workflow-sizes}. The largest experiment contains 50,000
tasks.

\begin{table}[t]
\centering
\caption{Workflow size categories used in the evaluation.}
\label{tab:workflow-sizes}
\begin{tabular}{lr}
\toprule
Category & Number of tasks \\
\midrule
Small  & 50--200 \\
Medium & 500--2,000 \\
Large  & 5,000--50,000 \\
\bottomrule
\end{tabular}
\end{table}

\subsection{Utilization and Credit-Stress Models}
\label{subsec:workload-models}

Four time-varying CPU-utilization models are used. The step model
changes abruptly from a low- to a high-utilization phase. The ramp
model increases utilization progressively over task execution. The
bursty model alternates low- and high-utilization phases, while the
trace-driven model superposes multiple temporal components to emulate
irregular production demand. Seeded perturbations and a bounded
co-location interference term are added to each model. Utilization is
clipped to the feasible interval $[0,1]$.

Three credit-stress regimes are evaluated. Under light stress,
aggregate demand exceeds baseline capacity infrequently. Moderate
stress is the default configuration and contains intermittent
above-baseline phases. Heavy stress introduces sustained high-demand
periods that rapidly deplete credits under credit-oblivious placement.
Credit is represented in vCPU-minutes, and the simulator applies the
credit evolution model in Eq.~(6) at every scheduling interval.

\subsection{Compared Scheduling Methods}
\label{subsec:baselines}

UBR-CA is compared with six references. HEFT performs
precedence-aware earliest-finish placement without modeling credits.
Kubernetes Resource Packing (KRP) performs deterministic CPU
bin-packing. Energy-Aware VM Consolidation (EA-VC) aggressively
consolidates tasks to reduce fragmentation. The Chance-Constrained
Scheduler (CCS) protects VM capacity using a fixed uncertainty bound
but does not model credit depletion. The Credit-Aware Reactive
Scheduler (CARS) reacts after a credit threshold is crossed. A
strict-feasibility branch-and-bound proxy provides a small-instance
reference for the robust placement rules.

The branch-and-bound result is a feasibility reference rather than a
claim of a globally optimal mixed-integer solution for large
workflows. Therefore, it is denoted ``B\&B proxy'' rather than
``MILP'' throughout this section. Two UBR-CA ablations are also
evaluated: one disables migration, and the other replaces Bayesian
posterior learning with deterministic utilization estimates.

\subsection{Implementation, Parameters, and Metrics}
\label{subsec:implementation}

The simulator and all schedulers were implemented in Java~21. The
implementation uses a discrete-event control loop with a default
scheduling interval of $\Delta t=300$~s, lookahead horizon $H=2$, and
risk tolerance $\epsilon=0.05$, corresponding to a 95\% one-sided
Bayesian bound. The prior variance is 0.04, the observation variance is
0.0225, the migration trigger is six credits, and the safe recovery
level is 12 credits. Table~\ref{tab:sensitivity-config} lists the
parameter values used in the sensitivity study.

\begin{table}[t]
\centering
\caption{Sensitivity-analysis configurations.}
\label{tab:sensitivity-config}
\begin{tabular}{ll}
\toprule
Parameter & Values \\
\midrule
$\Delta t$ & $\{60,300,900\}$ s \\
$H$ & $\{1,2,4\}$ \\
$\epsilon$ & $\{0.01,0.05,0.10\}$ \\
$\lambda$ & $\{10,100,1000\}$ \\
\bottomrule
\end{tabular}
\end{table}

Execution cost is the sum of VM price multiplied by active duration.
Deadline violation rate is the fraction of workflows that complete
after their deadlines, and total lateness is the sum of positive
completion-time excesses. A credit-exhaustion event is counted when a
VM enters an interval in which available credits cannot sustain its
above-baseline demand. Migration count records all completed
container migrations. Scheduler runtime measures CPU time spent in
placement, monitoring, and migration decisions; it excludes workload
execution time and artifact generation.

Every comparison is repeated using 20 seeds. A workflow set is
generated once per seed and cloned for every scheduler, producing
paired observations. Results are reported as mean $\pm$ sample
standard deviation. Cost comparisons use paired Student's $t$-tests,
Wilcoxon signed-rank tests, and paired Cohen's $d$ at a significance
level of $0.05$. All per-seed observations are retained in the
reproducibility package.

\subsection{Overall Performance}
\label{subsec:overall}

Table~\ref{tab:overall-performance} summarizes the default moderate
stress results. UBR-CA completed all workflows without deadline
violations and reduced mean credit-exhaustion events to
$2.10\pm1.21$. HEFT, KRP, EA-VC, and CCS each violated the deadlines
of approximately half of the workflows. CARS also met all deadlines,
but incurred $10.60\pm2.28$ credit-exhaustion events and
$95.45\pm16.66$ migrations. UBR-CA therefore reduced exhaustion
events by 80.19\% and migrations by 46.10\% relative to CARS.

The B\&B proxy and UBR-CA selected the same strict-feasibility
decisions in this configuration and consequently produced identical
execution metrics. This agreement is specific to the tested workload;
it does not establish global optimality for larger instances.

\begin{table*}[t]
\centering
\caption{Overall performance under moderate credit stress
(mean $\pm$ standard deviation over 20 paired seeds).}
\label{tab:overall-performance}
\scriptsize
\resizebox{\textwidth}{!}{%
\begin{tabular}{lrrrrr}
\toprule
Method & Cost (USD) & Deadline violations (\%) &
Credit exhaustions & Migrations & Runtime (s) \\
\midrule
HEFT        & $7.874\pm0.498$ & $50.00\pm0.00$ & $8.15\pm3.15$  & $0.00\pm0.00$   & $0.0228\pm0.0373$ \\
KRP         & $7.830\pm0.488$ & $50.00\pm0.00$ & $7.80\pm4.06$  & $0.00\pm0.00$   & $0.0157\pm0.0155$ \\
EA-VC       & $7.830\pm0.488$ & $50.00\pm0.00$ & $7.80\pm4.06$  & $0.00\pm0.00$   & $0.0157\pm0.0145$ \\
CCS         & $7.805\pm0.493$ & $50.00\pm0.00$ & $9.35\pm3.79$  & $0.00\pm0.00$   & $0.0163\pm0.0140$ \\
CARS        & $7.802\pm0.519$ & $0.00\pm0.00$  & $10.60\pm2.28$ & $95.45\pm16.66$ & $0.0149\pm0.0194$ \\
B\&B proxy  & $7.972\pm0.521$ & $0.00\pm0.00$  & $2.10\pm1.21$  & $51.45\pm9.23$  & $0.0198\pm0.0301$ \\
UBR-CA      & $7.972\pm0.521$ & $0.00\pm0.00$  & $2.10\pm1.21$  & $51.45\pm9.23$  & $0.0157\pm0.0201$ \\
\bottomrule
\end{tabular}}
\end{table*}

\subsection{Execution-Cost Analysis}
\label{subsec:cost}

Figure~\ref{fig:normalized-cost} reports cost normalized to HEFT.
All methods lie within a narrow cost band. UBR-CA incurred
$7.972\pm0.521$ USD, compared with $7.874\pm0.498$ USD for HEFT and
$7.802\pm0.519$ USD for CARS. Thus, the safety and deadline gains of
UBR-CA were obtained with a 1.25\% mean cost premium over HEFT and a
2.18\% premium over CARS in the synthetic default workload.

This result does not support an unconditional claim that UBR-CA
reduces cloud cost. Rather, it demonstrates a cost-safety trade-off:
the robust reserve and recovery migrations modestly increase billed
VM time while substantially reducing throttling risk and deadline
failure. A claim of cost reduction should be made only if a
trace-backed experiment using the supplied Alibaba importer produces
that outcome.

\begin{figure}[t]
\centering
\includegraphics[width=\columnwidth]{output/figures/figure_1_normalized_execution_cost.png}
\caption{Normalized workflow execution cost. Error bars show one
sample standard deviation across 20 paired seeds. The distributed
artifact is also available as an editable SVG and a 1600$\times$960
PNG; convert the SVG to PDF for vector inclusion if required.}
\label{fig:normalized-cost}
\end{figure}

\subsection{Workflow Deadline Compliance}
\label{subsec:deadline}

Table~\ref{tab:deadline-performance} shows that HEFT, KRP, EA-VC, and
CCS each produced a 50\% mean violation rate. Their lateness differed
substantially: HEFT incurred $47{,}845.3\pm11{,}952.5$~s, while KRP
and EA-VC incurred $134{,}653.5\pm11{,}272.0$~s. Capacity uncertainty
alone was insufficient: CCS still accumulated
$101{,}724.2\pm7{,}066.6$~s of lateness because it did not protect
future credit availability.

CARS and UBR-CA both completed all workflows before their deadlines.
The distinction between them is therefore not deadline compliance in
this configuration, but the number of unsafe credit states and
corrective migrations. Across all 20 pairs, UBR-CA improved the
deadline outcome relative to HEFT, KRP, EA-VC, and CCS. Because these
paired deadline differences have zero sample variance, a standardized
parametric effect size is ill-conditioned; the corresponding exact
two-sided sign-test probability is below $0.001$.

\begin{table}[t]
\centering
\caption{Workflow deadline performance
(mean $\pm$ standard deviation).}
\label{tab:deadline-performance}
\scriptsize
\begin{tabular}{lrr}
\toprule
Method & Violation rate (\%) & Total lateness (s) \\
\midrule
HEFT   & $50.00\pm0.00$ & $47{,}845.3\pm11{,}952.5$ \\
KRP    & $50.00\pm0.00$ & $134{,}653.5\pm11{,}272.0$ \\
EA-VC  & $50.00\pm0.00$ & $134{,}653.5\pm11{,}272.0$ \\
CCS    & $50.00\pm0.00$ & $101{,}724.2\pm7{,}066.6$ \\
CARS   & $0.00\pm0.00$  & $0.0\pm0.0$ \\
UBR-CA & $0.00\pm0.00$  & $0.0\pm0.0$ \\
\bottomrule
\end{tabular}
\end{table}

\subsection{Credit Dynamics and Stress Robustness}
\label{subsec:credits}

Figure~\ref{fig:credit-trajectory} shows the minimum, mean, and
maximum credit balances across active VMs for one representative
UBR-CA run. The lower trajectory approaches the migration threshold
during early high-demand phases, after which proactive redistribution
restores a positive reserve. The saw-tooth behavior of the upper
trajectory results from credit accumulation on lightly loaded VMs and
the subsequent retirement or reuse of those instances. Most
importantly, the lower trajectory does not exhibit sustained
zero-credit operation.

\begin{figure}[t]
\centering
\includegraphics[width=\columnwidth]{output/figures/figure_2_credit_trajectory.png}
\caption{Representative minimum, mean, and maximum CPU-credit
trajectories across active VMs under UBR-CA.}
\label{fig:credit-trajectory}
\end{figure}

Table~\ref{tab:credit-stress} and Fig.~\ref{fig:credit-stress} quantify
the effect of increasing stress. Under heavy stress, HEFT, KRP, and
CCS produced between 40.05 and 44.05 exhaustion events on average.
UBR-CA limited this value to $3.45\pm1.70$, corresponding to a 91.39\%
reduction relative to HEFT and a 91.85\% reduction relative to CCS.
Compared with CARS under heavy stress, UBR-CA reduced exhaustion
events by 69.20\%, while preserving zero deadline violations.
The heavy-stress safety gain carried a 4.47\% cost premium over HEFT
and a 0.57\% premium over CARS.

\begin{table*}[t]
\centering
\caption{Performance across credit-stress regimes
(mean $\pm$ standard deviation).}
\label{tab:credit-stress}
\scriptsize
\resizebox{\textwidth}{!}{%
\begin{tabular}{llrrr}
\toprule
Stress & Method & Cost (USD) & Violations (\%) & Exhaustions \\
\midrule
Light & HEFT   & $7.073\pm0.487$ & $50.00\pm0.00$ & $6.20\pm1.91$ \\
      & KRP    & $7.058\pm0.499$ & $50.00\pm0.00$ & $4.55\pm1.47$ \\
      & CCS    & $7.051\pm0.491$ & $50.00\pm0.00$ & $4.65\pm0.93$ \\
      & CARS   & $6.918\pm0.471$ & $0.00\pm0.00$  & $12.10\pm2.00$ \\
      & UBR-CA, no migration & $7.100\pm0.545$ & $0.00\pm0.00$ & $8.35\pm1.66$ \\
      & UBR-CA & $7.224\pm0.492$ & $0.00\pm0.00$  & $3.45\pm1.57$ \\
\midrule
Moderate & HEFT & $7.874\pm0.498$ & $50.00\pm0.00$ & $8.15\pm3.15$ \\
         & KRP  & $7.830\pm0.488$ & $50.00\pm0.00$ & $7.80\pm4.06$ \\
         & CCS  & $7.805\pm0.493$ & $50.00\pm0.00$ & $9.35\pm3.79$ \\
         & CARS & $7.802\pm0.519$ & $0.00\pm0.00$  & $10.60\pm2.28$ \\
         & UBR-CA, no migration & $7.855\pm0.553$ & $0.00\pm0.00$ & $13.20\pm2.26$ \\
         & UBR-CA & $7.972\pm0.521$ & $0.00\pm0.00$ & $2.10\pm1.21$ \\
\midrule
Heavy & HEFT   & $8.743\pm0.529$ & $50.00\pm0.00$ & $40.05\pm9.38$ \\
      & KRP    & $8.832\pm0.477$ & $55.00\pm6.28$ & $44.05\pm7.56$ \\
      & CCS    & $8.759\pm0.493$ & $50.00\pm0.00$ & $42.35\pm6.44$ \\
      & CARS   & $9.082\pm0.479$ & $0.00\pm0.00$  & $11.20\pm3.37$ \\
      & UBR-CA, no migration & $9.022\pm0.498$ & $0.00\pm0.00$ & $16.00\pm2.62$ \\
      & UBR-CA & $9.134\pm0.486$ & $0.00\pm0.00$  & $3.45\pm1.70$ \\
\bottomrule
\end{tabular}}
\end{table*}

\begin{figure}[t]
\centering
\includegraphics[width=\columnwidth]{output/figures/figure_9_credit_stress.png}
\caption{Mean credit-exhaustion events under light, moderate, and
heavy stress.}
\label{fig:credit-stress}
\end{figure}

\subsection{Ablation Study}
\label{subsec:ablation}

Table~\ref{tab:ablation} isolates the two major UBR-CA components.
Disabling migration increased credit-exhaustion events from
$2.10\pm1.21$ to $13.20\pm2.26$; equivalently, the full method reduced
exhaustion by 84.09\% relative to the no-migration variant, even though
all workflows still met their deadlines. This result shows that
robust initial placement alone cannot absorb all subsequent
utilization changes.

Replacing Bayesian learning with deterministic estimates reduced mean
cost from 7.972 to 7.818 USD but increased exhaustion events from 2.10
to 8.70 and migrations from 51.45 to 77.15. Hence Bayesian uncertainty
learning did not reduce monetary cost in this run; instead, it improved
credit safety and reduced corrective actions. Relative to the
deterministic variant, the full method reduced exhaustion by 75.86\%
and migration count by 33.31\%, at a 1.97\% cost premium. These
absolute metrics provide a more informative ablation than a
percentage deadline-reduction column because all three UBR-CA variants
had zero deadline violations.

\begin{table*}[t]
\centering
\caption{Ablation of Bayesian learning and migration
(mean $\pm$ standard deviation).}
\label{tab:ablation}
\scriptsize
\begin{tabular}{lrrrr}
\toprule
Variant & Cost (USD) & Violations (\%) & Exhaustions & Migrations \\
\midrule
UBR-CA without migration &
$7.855\pm0.553$ & $0.00\pm0.00$ & $13.20\pm2.26$ & $0.00\pm0.00$ \\
UBR-CA without Bayesian learning &
$7.818\pm0.480$ & $0.00\pm0.00$ & $8.70\pm2.34$ & $77.15\pm16.53$ \\
Full UBR-CA &
$7.972\pm0.521$ & $0.00\pm0.00$ & $2.10\pm1.21$ & $51.45\pm9.23$ \\
\bottomrule
\end{tabular}
\end{table*}

\subsection{Parameter Sensitivity}
\label{subsec:sensitivity}

Table~\ref{tab:sensitivity-results} presents the complete sensitivity
results. The scheduling interval exhibits the strongest trade-off.
At $\Delta t=60$~s, mean cost was lowest, but the configured
short-horizon reserve admitted $18.80\pm4.63$ exhaustion events. A
300-s interval reduced this value to $2.10\pm1.21$. At 900~s,
exhaustion was nearly eliminated, but delayed decision epochs raised
cost to $8.931\pm0.524$, 12.03\% above the 300-s setting.

Increasing $H$ strengthened credit protection. Moving from $H=1$ to
$H=2$ reduced exhaustion from 8.85 to 2.10 events, while $H=4$
eliminated exhaustion at a 1.47\% cost increase over $H=2$. The risk
parameter behaved as expected: $\epsilon=0.01$ produced fewer
exhaustions than $\epsilon=0.10$ (1.95 versus 2.85), but cost 1.33\%
more. Varying $\lambda$ did not change the result because every
UBR-CA workflow met its deadline; consequently, the lateness term was
inactive throughout these runs.

\begin{table*}[t]
\centering
\caption{UBR-CA parameter sensitivity
(mean $\pm$ standard deviation).}
\label{tab:sensitivity-results}
\scriptsize
\resizebox{\textwidth}{!}{%
\begin{tabular}{llrrr}
\toprule
Parameter & Value & Cost (USD) & Violations (\%) & Exhaustions \\
\midrule
$\Delta t$ (s) & 60  & $7.820\pm0.445$ & $0.00\pm0.00$ & $18.80\pm4.63$ \\
               & 300 & $7.972\pm0.521$ & $0.00\pm0.00$ & $2.10\pm1.21$ \\
               & 900 & $8.931\pm0.524$ & $0.00\pm0.00$ & $0.05\pm0.22$ \\
\midrule
$H$ & 1 & $7.958\pm0.540$ & $0.00\pm0.00$ & $8.85\pm2.58$ \\
    & 2 & $7.972\pm0.521$ & $0.00\pm0.00$ & $2.10\pm1.21$ \\
    & 4 & $8.089\pm0.549$ & $0.00\pm0.00$ & $0.00\pm0.00$ \\
\midrule
$\epsilon$ & 0.01 & $8.018\pm0.523$ & $0.00\pm0.00$ & $1.95\pm1.47$ \\
           & 0.05 & $7.972\pm0.521$ & $0.00\pm0.00$ & $2.10\pm1.21$ \\
           & 0.10 & $7.913\pm0.573$ & $0.00\pm0.00$ & $2.85\pm1.18$ \\
\midrule
$\lambda$ & 10   & $7.972\pm0.521$ & $0.00\pm0.00$ & $2.10\pm1.21$ \\
          & 100  & $7.972\pm0.521$ & $0.00\pm0.00$ & $2.10\pm1.21$ \\
          & 1000 & $7.972\pm0.521$ & $0.00\pm0.00$ & $2.10\pm1.21$ \\
\bottomrule
\end{tabular}}
\end{table*}

Figure~\ref{fig:sensitivity-panels} visualizes these four parameter
sweeps. Because cost, violation rate, and exhaustion count use
different units, the table should be used for precise comparison; the
plots show trends rather than a common-scale objective.

\begin{figure*}[t]
\centering
\begin{minipage}{0.49\textwidth}
\centering
\includegraphics[width=\linewidth]{output/figures/figure_4_sensitivity_interval_seconds.png}\\[-1mm]
\small (a) Scheduling interval
\end{minipage}
\hfill
\begin{minipage}{0.49\textwidth}
\centering
\includegraphics[width=\linewidth]{output/figures/figure_5_sensitivity_lookahead_H.png}\\[-1mm]
\small (b) Lookahead horizon
\end{minipage}

\vspace{1mm}
\begin{minipage}{0.49\textwidth}
\centering
\includegraphics[width=\linewidth]{output/figures/figure_6_sensitivity_epsilon.png}\\[-1mm]
\small (c) Risk tolerance
\end{minipage}
\hfill
\begin{minipage}{0.49\textwidth}
\centering
\includegraphics[width=\linewidth]{output/figures/figure_7_sensitivity_lateness_lambda.png}\\[-1mm]
\small (d) Lateness penalty
\end{minipage}
\caption{Sensitivity of UBR-CA to its main control parameters.}
\label{fig:sensitivity-panels}
\end{figure*}

\subsection{Bayesian Uncertainty Convergence}
\label{subsec:posterior-convergence}

Figure~\ref{fig:posterior-convergence} illustrates the analytical
posterior-variance contraction used by the scheduler. Both the
posterior standard deviation and the 95\% uncertainty half-width
decrease monotonically as utilization observations accumulate. The
result explains why UBR-CA can reserve additional capacity during a
task's poorly observed initial phase and reduce estimation
conservatism later. The ablation results show that this calibration
primarily improves credit safety and migration efficiency in the
present workload rather than reducing cost.

\begin{figure}[t]
\centering
\includegraphics[width=\columnwidth]{output/figures/figure_8_bayesian_convergence.png}
\caption{Posterior standard deviation and 95\% uncertainty half-width
as the number of task-utilization observations increases.}
\label{fig:posterior-convergence}
\end{figure}

\subsection{Scalability}
\label{subsec:scalability}

Figure~\ref{fig:scalability} reports the scheduler CPU time as workflow
size increases. UBR-CA required 0.000082~s for 50 tasks,
0.001348~s for 200 tasks, 0.002285~s for 500 tasks, 0.011746~s for
2,000 tasks, 0.045420~s for 5,000 tasks, 0.208469~s for 10,000 tasks,
and 2.033631~s for 50,000 tasks. The measured growth is consistent
with the placement and migration complexity in Eq.~(84), while
remaining well below the 300-s default scheduling interval at every
tested scale. These values measure scheduler computation only and
should not be interpreted as end-to-end workflow simulation time.

\begin{figure}[t]
\centering
\includegraphics[width=\columnwidth]{output/figures/figure_3_scheduler_scalability.png}
\caption{UBR-CA scheduler CPU time versus workflow size. The horizontal
axis uses a logarithmic scale.}
\label{fig:scalability}
\end{figure}

\subsection{Statistical Analysis}
\label{subsec:statistics}

Table~\ref{tab:cost-significance} reports paired tests for execution
cost. The comparison with HEFT is not significant at the 0.05 level
($p=0.0610$). The comparisons with KRP, CCS, and CARS are significant,
but the positive paired Cohen's $d$ values indicate that the
difference is a UBR-CA cost premium, not a cost saving. This direction
must be retained when interpreting the tests.

The zero-variance deadline outcomes make a conventional standardized
effect size inappropriate. Descriptively, UBR-CA improved every paired
deadline outcome relative to HEFT, KRP, EA-VC, and CCS, while matching
CARS. Credit-exhaustion and migration results are reported with their
full per-seed observations in the reproducibility package so that
metric-specific nonparametric analyses can be applied without
discarding ties.

\begin{table}[t]
\centering
\caption{Paired execution-cost comparisons between UBR-CA and the
principal baselines. Positive $d$ denotes higher UBR-CA cost.}
\label{tab:cost-significance}
\scriptsize
\begin{tabular}{lrrr}
\toprule
Comparison & $t$-test $p$ & Wilcoxon $p$ & Cohen's $d$ \\
\midrule
UBR-CA vs.\ HEFT & 0.0610 & 0.0910 & 0.445 \\
UBR-CA vs.\ KRP  & 0.0047 & 0.0080 & 0.715 \\
UBR-CA vs.\ CCS  & 0.0014 & 0.0054 & 0.835 \\
UBR-CA vs.\ CARS & $<0.001$ & $<0.001$ & 1.246 \\
\bottomrule
\end{tabular}
\end{table}

\subsection{Discussion and Validity}
\label{subsec:discussion}

The experiments support three main conclusions. First, explicitly
protecting future credit availability is necessary for reliable
workflow execution. Capacity-aware methods that omit credit dynamics
produced substantial lateness, and the advantage of UBR-CA increased
sharply under heavy stress. Second, Bayesian learning and migration
are complementary. Migration supplies the corrective action required
when demand changes after placement, while Bayesian bounds reduce the
number of unsafe states and unnecessary corrective actions. Third,
these reliability gains are not free in the present configuration.
UBR-CA exchanged a small increase in monetary cost for zero deadline
violations, fewer exhaustion events, and fewer migrations than the
reactive baseline.

The results also delimit the claims that can be made. The reported
workloads are benchmark-shaped and synthetic; trace-backed Alibaba
claims require a separate run using the provided importer. CPU is the
only constrained resource, so the evaluation does not capture memory,
storage, image-pull delay, or network contention. The B\&B method is a
strict-feasibility proxy rather than a global large-instance optimum.
Finally, scheduler runtime is measured within the Java implementation
and is hardware dependent. These limitations do not change the
observed relative credit-safety behavior, but they should accompany
any generalization to a specific cloud provider or production cluster.

Overall, the evidence supports UBR-CA as a deadline- and credit-safety
mechanism with bounded online overhead. For the reported synthetic
workloads, its principal benefit is reliability rather than lower
cloud expenditure.
